\chapter{Scotoma}

\section*{Definition}
A scotoma is an area of partial or complete visual field loss surrounded by normal or relatively normal vision, quantified on perimetry as a depression of sensitivity that may be central, paracentral, arcuate, or altitudinal in location.

\section*{What Happens in Your Body}
Scotomas arise from damage anywhere in the visual pathway. Retinal causes (macular disease) produce central scotomas; optic nerve disease produces arcuate or altitudinal patterns; post-chiasmal lesions produce homonymous congruous defects. The underlying mechanism is interruption of neural signal transmission at the affected level.

\section*{Causes}
\begin{itemize}
  \item \textbf{Macular:} Age-related macular degeneration, macular hole, central serous chorioretinopathy
  \item \textbf{Optic nerve:} Glaucoma (arcuate scotoma), optic neuritis, anterior ischemic optic neuropathy
  \item \textbf{Retinal vascular:} Branch retinal artery or vein occlusion
  \item \textbf{Retinal detachment:} Large peripheral scotoma corresponding to detached retina
  \item \textbf{Brain:} Occipital lobe stroke or tumor (homonymous hemianopia)
  \item \textbf{Migraine:} Transient scintillating scotoma (fortification spectra) during aura
\end{itemize}

\section*{When to See a Doctor}
See your doctor if:
\begin{itemize}
  \item Sudden onset of any blind spot in the visual field
  \item Progressive enlargement of an existing scotoma
  \item Scotoma associated with flashing lights or floaters
  \item Transient scotoma consistent with amaurosis fugax or migraine (evaluate to distinguish)
  \item Bilateral homonymous defect concerning for intracranial pathology
\end{itemize}

\section*{Related Symptoms}
\begin{itemize}
  \item Visual field defect (altitudinal, arcuate, central, or homonymous)
  \item Decreased visual acuity (if scotoma involves the fovea)
  \item Metamorphopsia (distortion of straight lines)
  \item Photopsia (if retinal irritation)
  \item Color desaturation (particularly red in optic nerve disease)
\end{itemize}
\textbf{Important:} Scotomas are often unnoticed by you until formally tested (visual completion phenomenon); people with significant field loss may still read 20/20 on the Snellen chart.

\section*{Self-Care / Home Management}
\begin{itemize}
  \item Amsler grid self-testing at home for monitoring central scotoma progression
  \item Scanning techniques: training the eyes to move systematically to compensate for field loss
  \item Eccentric viewing training for central scotoma (using preferred retinal locus)
  \item Improved lighting and magnification for reading with central field loss
  \item Driving cessation if the scotoma encroaches on the central 20\textdegree{} of vision
\end{itemize}

\section*{How Your Doctor Will Evaluate You}
\begin{itemize}
  \item Amsler grid testing for central scotoma detection and monitoring
  \item Confrontation visual field testing at bedside
  \item Automated static perimetry (Humphrey 24-2 or 30-2 SITA) for formal field mapping
  \item Optical coherence tomography (OCT) of macula and optic nerve head
  \item MRI brain with contrast if homonymous or bitemporal defects suggest intracranial lesion
  \item Carotid Doppler or echocardiography if embolic source suspected (altitudinal defect)
\end{itemize}

\section*{Treatment Options}
\begin{itemize}
  \item \textbf{Macular disease:} Anti-VEGF injections for neovascular age-related macular degeneration; observation for macular hole (with vitrectomy if progressive)
  \item \textbf{Glaucoma:} IOP-lowering therapy to prevent scotoma progression (existing defects are irreversible)
  \item \textbf{Optic neuritis:} IV steroids to hasten recovery; long-term immunomodulation in multiple sclerosis
  \item \textbf{Retinal detachment:} Surgical reattachment to restore peripheral field
  \item \textbf{Stroke:} Low-vision rehabilitation with prismatic field-expansion devices
  \item \textbf{Migraine:} Abortive triptan therapy and prophylactic agents (beta-blockers, topiramate)
\end{itemize}

\section*{References}
\begin{thebibliography}{99}
\bibitem{scotoma:ref1} Bowd C, Zangwill LM, Berry CC, et al. "Detecting early glaucoma by assessment of retinal nerve fiber layer thickness and visual function." \textit{Investigative Ophthalmology \& Visual Science}, 2001;42(9):1993--2003.
\bibitem{scotoma:ref2} Balcer LJ, Miller DH, Reingold SC, et al. "Vision and vision-related outcome measures in multiple sclerosis." \textit{Brain}, 2010;133(3):700--711.
\bibitem{scotoma:ref3} Hayreh SS, Zimmerman MB, Podhajsky PA, et al. "Nonarteritic anterior ischemic optic neuropathy: Natural history of visual outcome." \textit{Ophthalmology}, 2008;115(11):1936--1945.
\bibitem{scotoma:ref4} Wong TY, Chakravarthy U, Klein R, et al. "The natural history and prognosis of neovascular age-related macular degeneration: A systematic review." \textit{Ophthalmology}, 2008;115(1):116--126.
\bibitem{scotoma:ref5} Gazzard G, Konstantakopoulou E, Garway-Heath D, et al. "Selective laser trabeculoplasty versus eye drops for first-line glaucoma treatment (LiGHT trial)." \textit{The Lancet}, 2019;393(10180):1505--1516.
\bibitem{scotoma:ref6} Scholl HPN, Strauss RW, Singh MS, et al. "Emerging therapies for inherited retinal degeneration." \textit{Science Translational Medicine}, 2016;8(368):368rv6.
\end{thebibliography}
